    
\documentclass[11pt]{article}
\usepackage{times}
    \usepackage{fullpage}
    
    \title{Analysis of gender stereotypes in personality assessments}
    \author{Sofia Karamichali - 2570646K}

    \begin{document}
    \maketitle
    
    
    
\section{Status report}

\subsection{Proposal}\label{proposal}

\subsubsection{Motivation}\label{motivation}

\emph{This project is motivated by the importance of analysis of datasets via scientific and
statistical methodologies and the need for the identification of biases within them. Thus,
informed use of data can be achieved for purposes such as machine learning and the
development of effective systems without inappropriate biases. The connection between
social perception and the emerging stereotypes further motivates the importance of this
project.}

\subsubsection{Aims}\label{aims}

\emph{This project aims to determine gender biases present in personality judgements based on speech.
The data used comes from a dataset of audio clips judged by listeners in terms of
personality traits. The analysis primarily focuses on the distinction between the speakers'
gender and how this impacted the judges’ perceptions highlighting the presence of any bias
in the collected data. The project involves various methodologies used to
determine the stereotypes identified in the dataset.}

\subsection{Progress}\label{progress}
\begin{itemize}
    \item Background research done on Big Five personality traits and personality
    perception.
    \item Background research done on statistical tests and hypothesis testing methodologies.
    \item Analysis of the overall dataset agreement using statistical measures to establish
    dataset reliability done.
    \item Analysis of personality trait data to define patterns of gender bias.
    \item Background research conducted on pattern recognition, machine learning and use
    of the Scikit-Learn python library.
    \item Implementation of basic pattern recognition to determine distinction between
    genders by applying data to test and train models done.
    \item Statistical methodologies applied to define the credibility of the models’ outcomes.
\end{itemize}


\subsection{Problems and risks}\label{problems-and-risks}

\subsubsection{Problems}\label{problems}

\begin{itemize}
    \item Delayed assignment of the project
    \item Lack of background knowledge of the statistical methods
    \item Lack of background knowledge of machine learning and pattern recognition
    \item Large imbalance in dataset gender split – Used appropriate tests and methods to
    handle the imbalance.
\end{itemize}

\subsubsection{Risks}\label{risks}
Lack of other biases in the dataset \& focus of the project on gender biases - might not allow for experiemntation of other biases. 
\\Solution - Adopt extended analysis \& experimentation of gender biases

\subsection{Plan}\label{plan}

\emph{Semester 2}
\begin{itemize}
    \item Week 1-2: Complete further research \& identify further experiments to be
    completed.
    \\Deliverable: Literature review draft \& more specific plan for upcoming weeks
    \item Week 3-5: Implement any further experiments \& visualizations \& collect evaluated
    results.
    \\Deliverable: Completion of experiments \& detailed findings
    \item Week 6: Final implementation of all experiments
    \\Deliverable: polished software ready
    \item Week 7-8: Write up.
    First draft.
    \item Week 9-10: Write up.
    \\Deliverable: Complete final draft.
\end{itemize}

    
\subsection{Ethics and data}\label{ethics}
This project does not involve human subjects or data. No approval required.



\end{document}
